\section{Algorithme CLAPOTI}

Être polynomial en $\log(|\Delta|)$ $=$ éviter de calculer dans le groupe de classe. Nouvelle méthode avec Kani

A citer : Some applications of higher dimensional isogenies to elliptic curves, Damiens Robert

\subsection{Lemme de Kani}

\begin{Def}
	Un diamant d'isogénie $\left( E,f_{1}, f'_{1}, f_2, f'_2 \right)$  entre courbe elliptiques est un diagramme commutatif de la forme 
	\[\begin{tikzcd}
		E & {E_1} \\
		{E_2} & {E'}
		\arrow["{f_1}", from=1-1, to=1-2]
		\arrow["{f_2}"', from=1-1, to=2-1]
		\arrow["{f'_1}", from=1-2, to=2-2]
		\arrow["{f'_2}"', from=2-1, to=2-2]
	\end{tikzcd}\]
	avec $E$, $E_1$, $E_2$, $E'$ des courbes elliptiques, $f_1$, $f'_2$ deux $d_1$-isogénies, $f_2$, $f'_1$ deux $d_2$-isogénies, et $f = f_1\circ f'_1 = f_2\circ f'_2$ une $d_1d_2$-isogénie. 
\end{Def}

\begin{Prop}
	Soit $\left( E,f_{1}, f'_{1}, f_2, f'_2 \right)$ un diamant d'isogénie. On considère
	\begin{equation*}
		F : E\times E' \rightarrow E_1\times E_2\;,\;(P,Q)\mapsto \left( f_1(P) + \hat{f}_1(Q), \hat{f}_2(Q) - f_2(P) \right), 
	\end{equation*}
	qui est une isogénie entre produit de courbes elliptiques. Alors
	\begin{equation*}
		\ker(F) = \left\lbrace  \right\rbrace 
	\end{equation*}
\end{Prop}




\subsection{Calcul d'isogénies entre produits de courbes elliptiques}

\subsection{N-torsion accessible}

